% Generated draft: requires mathematical and source review.
\documentclass[11pt]{article}
\usepackage[margin=25mm]{geometry}
\usepackage{fontspec}
\setmainfont{Latin Modern Roman}
\usepackage{amsmath,amssymb,mathtools,unicode-math}
\setmathfont{Latin Modern Math}
\usepackage{enumitem,needspace,xurl,hyperref}
\setlength{\parindent}{0pt}
\setlength{\parskip}{5pt}
\setlength{\emergencystretch}{3em}
\newcommand{\N}{\mathbb{N}}
\newcommand{\Z}{\mathbb{Z}}
\newcommand{\Q}{\mathbb{Q}}
\newcommand{\R}{\mathbb{R}}
\newcommand{\C}{\mathbb{C}}
\newcommand{\F}{\mathbb{F}}
\newcommand{\A}{\mathbb{A}}
\newcommand{\PP}{\mathbb{P}}
\newcommand{\E}{\mathbb{E}}
\newcommand{\eps}{\varepsilon}
\newcommand{\dd}{\,\mathrm{d}}
\begin{document}
\section*{UnsolvedMath ID: 30002786}
\textbf{Source problem code:} catalogue entry
\subsection{Definitions and mathematical statement}
Let \(M=S(S^3)\) denote the total space of the real spinor bundle over \(S^3\), a seven-dimensional manifold. Here \(S(S^3)\) is interpreted as the total bundle, not its unit sphere bundle, consistently with the supplied description of a holonomy-\(G_2\) metric.

For clarity, a standard meaning of an asymptotically conical (AC) metric \(g\) on \(M\) is the following: there exist a compact subset \(K\subset M\), a compact six-dimensional manifold \(\Sigma\) with a Riemannian metric \(h\), a number \(R>0\), and a diffeomorphism
\[
\Phi:(R,\infty)\times\Sigma\longrightarrow M\setminus K
\]
such that, with \(g_C=\mathrm{d}r^2+r^2h\), there is a rate \(\nu<0\) for which
\[
\left|\nabla_{g_C}^{\,j}(\Phi^*g-g_C)\right|_{g_C}
=O(r^{\nu-j})\qquad(r\to\infty)
\]
for every integer \(j\geq0\). This states a standard smooth AC convention; the supplied record specifies neither a decay rate nor whether the asymptotic cone is prescribed. The phrase holonomy-\(G_2\) is read as \(\operatorname{Hol}(g)=G_2\); whether holonomy merely contained in \(G_2\) was intended requires confirmation.

\textbf{Source question.} Is the asymptotically conical holonomy-\(G_2\) metric on \(M\), identified in the supplied title and background with the Bryant--Salamon metric, unique?

\textbf{Unresolved precision of uniqueness.} Let \(g_{\mathrm{BS}}\) denote the Bryant--Salamon metric referred to by the record. One possible precise interpretation asks whether every complete AC metric \(g\) on \(M\) with \(\operatorname{Hol}(g)=G_2\) satisfies
\[
 g=c^2F^*g_{\mathrm{BS}}
 \quad\text{for some }c>0\text{ and diffeomorphism }F:M\longrightarrow M.
\]
This is a conditional formulation for review, not an assertion that the source specifies these hypotheses or this equivalence relation. The record does not explicitly require completeness, identify the allowed asymptotic cones or rates, fix a scale, or specify which diffeomorphisms are allowed. Literal uniqueness cannot distinguish a metric from its constant rescalings and diffeomorphic pullbacks without additional normalization or an equivalence relation. The background mentions cohomogeneity-one classification, but the question itself imposes no symmetry assumption; restriction to that class is therefore unresolved and must not be inferred.
\subsection{Short English statement}
The supplied question asks whether the Bryant--Salamon asymptotically conical metric with holonomy \(G_2\) on the total spinor bundle of \(S^3\) is unique. A precise uniqueness claim requires specifying completeness, the asymptotic conditions, and equivalence under diffeomorphisms and rescaling. The conditional formulation above makes one possible interpretation explicit and requires human mathematical review.
\subsection{Sources}
\begin{itemize}
\item Supplied local UnsolvedMath registry record: permanent ID \(30002786\), source problem code \texttt{catalogue entry}, in the Oberwolfach Reports Open Problems collection. The record gives the proposed year as 2015 and supplies no individual proposer.
\item \emph{Mini-Workshop: Singularities in \(G_2\)-geometry} (2015), \emph{Oberwolfach Reports}, pp.\ 449--488, DOI \texttt{10.4171/owr/2015/8}. Supplied source report and literature reference labelled ``original OWR report'': \url{https://doi.org/10.4171/owr/2015/8}.
\item Supplied literature reference labelled ``Foscolo, Haskins, and Nordstrom, cohomogeneity-one G2 manifolds'': \url{https://arxiv.org/abs/1805.02612}.
\item Literature assessment dated \texttt{2026-08-21}: the supplied curation labels the problem open and states that the known Bryant--Salamon AC \(G_2\) metrics on the spin bundle of \(S^3\) are locally unique, while global uniqueness remains tied to an unresolved classification of cohomogeneity-one \(G_2\) manifolds. This is supplied dated curation, not an independently verified literature assessment; the record itself requests independent verification of time-sensitive status and cited claims.
\item The registry labels its statement \texttt{corrected\_verified} and reports manual semantic, typographic, and LaTeX review of the cleanup. These are supplied editorial assertions, not independent verification of the intended uniqueness formulation or mathematical correctness.
\end{itemize}
\end{document}
